%----------------------------------------------------------------------------
\chapter{Background and theoretical foundations}
\label{chap:Background}
%----------------------------------------------------------------------------

\begin{itemize}
    \item This chapter establishes the mathematical foundations for the rest of the thesis.
    \item The central object is the variational form of implicit Euler, which unifies XPBD and VBD as two approximate minimisers of the same energy function, and therefore justifies a shared Newton reference as ground truth.
    \item Theory from \cite{li2026physics}. It contains detaild descriptions and case-studies about the theory and algorithms of most areas of physics-based simulations.
\end{itemize}

%----------------------------------------------------------------------------
\section{Time-stepping schemes for physical simulation}
\label{sec:Integration}
%----------------------------------------------------------------------------
\begin{itemize}
    \item Real-time physics reduces to advancing a state $(\mathbf{x}, \mathbf{v})$ --- positions and velocities of $n$ particles --- forward in time.
    \item In modern physics engines three families of approaches can be identified by splitting them based on how they resolve constraints and calculate motion.
\end{itemize}

\paragraph{Force-based, Newtonian.}
\begin{itemize}
    \item Forces $\mathbf{f}(\mathbf{x}, \mathbf{v}, t)$ produce accelerations via Newton's second law $f = M \dfrac{dv}{dt}$ and $v = \dfrac{dx}{dt}$.
    \item We get $\ddot{\mathbf{x}} = \mathbf{M}^{-1} \mathbf{f}$.
    \item Where $\mathbf{M}$ is the mass matrix.
    \item The most intuitive approach.
    \item This is mostly obsolete for real-time applications, but is still used for scenarios where high precision is required (engineering problems, offline VFX).
\end{itemize}

\paragraph{Impulse-based.}
\begin{itemize}
    \item Velocities are manipulated directly via collision or constraint impulses.
    \item Popular for rigid bodies, as it handles collisions as instantaneous changes.
    \item Not pursued here, because instantaneous change is not suitable for deformable bodies.
\end{itemize}

\paragraph{Position-based.}
\begin{itemize}
    \item Resolves consstraints by directly updating positions, without explicitly calculating forces or accelerations.
    \item PBD and XPBD belong in this cathegory.
    \item VBD is a hybrid: it is mathematically a position-based solver, as it directly updates positions, but does it by solving a local optimization problem that mimics a forced-based solver, so it is also a force-based solver in spirit.
\end{itemize}

\subsection{Explicit Time Integration (Forward Euler) and why it fails}
\label{ssec:ExplicitEuler}
It's explicit because it's a direct method to calculate $\mathbf{x}_{t+h}$ and $\mathbf{v}_{t+h}$ by substituting known values into simple formulas.
The explicit Euler scheme is
\begin{align}
    \mathbf{v}_{t+h} &= \mathbf{v}_t + h\,\mathbf{M}^{-1}\mathbf{f}(\mathbf{x}_t), \\
    \mathbf{x}_{t+h} &= \mathbf{x}_t + h\,\mathbf{v}_{t+h}.
\end{align}
\begin{itemize}
    \item Where $\mathbf{v_t}$ and $\mathbf{x}_t$ are the velocity and position at time $t$, and $h$ is the time step.
    \item This method is considered unconditionally unstable, meaning that irrespective of $h$ the simulation will eventually explode for equations with nonconstant $\mathbf{f}$.
\end{itemize}

\subsection{Implicit Time Integration (Backward Euler)}
\label{ssec:ImplicitEuler}
In contrast, implicit time integration requires solving a system of equations to determine the values of $\mathbf{x}_{t+h}$ and $\mathbf{v}_{t+h}$.
The (backward) implicit Euler scheme evaluates forces at the unknown
end-of-step state:
\begin{align}
    \mathbf{v}_{t+h} &= \mathbf{v}_t + h\,\mathbf{M}^{-1}\mathbf{f}(\mathbf{x}_{t+h}), \\
    \mathbf{x}_{t+h} &= \mathbf{x}_t + h\,\mathbf{v}_{t+h}.
\end{align}
\begin{itemize}
    \item Note that the calculation of $\mathbf{v}_{t+h}$ depends on $\mathbf{x}_{t+h}$, and $\mathbf{x}_{t+h}$ depends on $\mathbf{v}_{t+h}$.
    \item This gets resolved by substituting the$\mathbf{v}_{t+h} = (\mathbf{x}_{t+h} - \mathbf{x}_t)/h$ estimate into the position update equation.
    \item Unconditionally stable for linear systems.
    \item If the approximation is far from the actual solution, the method might not converge.
    \item Cost: a non-linear root-finding problem each step.
\end{itemize}

%----------------------------------------------------------------------------
\section{Implicit Euler as energy minimisation}
\label{sec:Variational}
%----------------------------------------------------------------------------
Eliminating the velocity equation gives the well-known variational form (also called minimization or optimization form) of
implicit Euler: $x$ is the minimiser of the energy function $G$:
$x_{t+h} = \argmin_{\mathbf{x}} G(\mathbf{x})$
\begin{equation}
    \label{eq:G}
    G(\mathbf{x}) \;=\; \tfrac{1}{2 h^{2}} \,\|\mathbf{x} - \mathbf{y}\|_{\mathbf{M}}^{2}
                       \;+\; E(\mathbf{x}),
\end{equation}
where
\begin{equation*}
    \mathbf{y} \;=\; \mathbf{x}_{t} + h\,\mathbf{v}_{t} + h^{2}\,\mathbf{M}^{-1}\mathbf{f}_{\mathrm{ext}}
\end{equation*}
\begin{itemize}
    \item $\mathbf{y}$ is the \emph{inertial target} (the position reached under external forces alone).
    \item $\|\cdot\|_{\mathbf{M}}$ is the mass-weighted length; $E(\mathbf{x})$ is the total potential energy evaluated at $\mathbf{x}$.
    \item Equation~\eqref{eq:G} is central: \textbf{both XPBD and VBD are approximate minimisers of the same $G(\mathbf{x})$}.
    \item They differ in \emph{which unknowns} they optimise over. VBD minimises $G$ directly over the positions $\mathbf{x}$ --- the \emph{primal} variables --- by block-coordinate descent. XPBD instead rewrites the elastic energy as a set of compliant constraints and solves for the Lagrange multipliers (the constraint forces) that enforce them --- the \emph{dual} variables. This is precisely what makes one a primal and the other a dual method; \autoref{sec:PrimalDual} develops the distinction and what it predicts.
    \item This shared target licenses a Newton solve of \eqref{eq:G} to serve as unambiguous ground truth in \autoref{chap:Experiments}.
\end{itemize}

%----------------------------------------------------------------------------
\section{Mass-spring systems}
\label{sec:MassSpring}
%----------------------------------------------------------------------------
\begin{itemize}
    \item I use a mass-spring discretisation of cloth: vertices carry point masses, edges of the cloth mesh carry springs.
    \item For a spring between particles $i$ and $j$ with rest length $\ell_0$ and stiffness $k$, the constraint function and energy are:
\end{itemize}
\begin{align}
    C(\mathbf{x}_i, \mathbf{x}_j) &= \|\mathbf{x}_i - \mathbf{x}_j\| - \ell_0, \\
    E(\mathbf{x}_i, \mathbf{x}_j) &= \tfrac{1}{2}\, k\, C^{2}.
\end{align}

I will introduce the exact mass-spring structure used to represent cloth in \autoref{ssec:ClothModel}.

\subsection{Limitations of mass-spring}
\label{ssec:MSSLimits}
\begin{itemize}
    \item Mass-spring is cheap and easy to differentiate, but is well known to be \emph{simulation-mesh-dependent}: effective material parameters, like macroscopic stiffness, depend on the discretisation.
    \item Therefore conclusions about absolute material parameters are not transferable to FEM-based models.
    \item Conclusions about \emph{relative} solver behaviour --- my subject --- are transferable.
\end{itemize}

%----------------------------------------------------------------------------
\section{Solver structure: Gauss--Seidel vs.\ Jacobi}
\label{sec:GSJacobi}
%----------------------------------------------------------------------------
\begin{itemize}
    \item Both XPBD's constraint sweep and VBD's vertex-block sweep are Gauss--Seidel-style iterations: each local update uses the most recent neighbour values, not start-of-iteration values.
    \item A Jacobi variant updates all blocks from the start-of-iteration state in parallel: slower convergence in serial, but maps better to the GPU.
    \item Consequence (revisited in \autoref{sec:ThreatsValidity}): both solvers are \textbf{order-dependent}.
    \begin{itemize}
        \item Constraint-traversal order (XPBD) and vertex-traversal order (VBD) both affect convergence rate and, in some (but relatively rare) cases, the converged solution.
        \item My experiments fix a deterministic order for reproducibility, I do not use randomisation, as none of the canonical solver implementations do.
    \end{itemize}
\end{itemize}

%----------------------------------------------------------------------------
\section{The primal/dual distinction}
\label{sec:PrimalDual}
%----------------------------------------------------------------------------
\begin{itemize}
    \item The unifying framing is the primal/dual perspective due to Macklin et~al.~\cite{10.1111/cgf.14104}.
    \item Minimising $G(\mathbf{x})$ in \eqref{eq:G} can be approached two ways:
\end{itemize}

\paragraph{Primal.}
\begin{itemize}
    \item Directly minimise over positions $\mathbf{x}$.
    \item A vertex-by-vertex block-coordinate descent on $G$ is exactly VBD.
\end{itemize}

\paragraph{Dual.}
\begin{itemize}
    \item Reformulate the energy through compliant constraints (each spring is a soft constraint with compliance $\alpha = 1/k$) and solve for Lagrange multipliers $\boldsymbol{\lambda}$ on each constraint.
    \item Introducing the multipliers turns the pure minimisation over $\mathbf{x}$ into a coupled problem in both $\mathbf{x}$ and $\boldsymbol{\lambda}$: a \emph{saddle-point} (or KKT) system, so called because its solution minimises the energy in the position variables while maximising it in the multipliers, rather than minimising in all unknowns at once.
    \item XPBD solves this system with a Gauss--Seidel sweep, updating one constraint's multiplier --- and the positions it touches --- at a time.
\end{itemize}

This duality is the lens is the motivation in the thesis for choosing the stress tests in \autoref{ssec:ChainWithHeavyWeight} and \autoref{ssec:StiffnessRatioChain}.
It also shows us where the two
approaches naturally differ:
\begin{itemize}
    \item Dual variables couple constraints that share a particle, so mass disparities at a shared particle are a scenario where XPBD's per-constraint updates may propagate slowly.
    \item The primal Hessian block at a vertex collects all spring stiffnesses meeting there, so widely varying stiffnesses produce an ill-conditioned block --- a scenario where VBD's per-vertex update may struggle.
\end{itemize}